% !PRINTED 673
\weberpage{673}

Ist $\mathfrak{P}$ irgend ein Punkt und $\varrho$ eine Funktion, die in diesem
Punkt die Ordnung 1 hat, so setzen wir nach § 185

\weq{2}{\begin{aligned}
\eta_1 &= e_1\varrho^m + \sigma_1\varrho^{m+1},\\
\eta_2 &= e_2\varrho^m + \sigma_2\varrho^{m+1},\\
&\text{\quad$\cdot\ \cdot\ \cdot\ \cdot\ \cdot\ \cdot\ \cdot\ \cdot\ \cdot\ \cdot$}\\
\eta_s &= e_s\varrho^m + \sigma_s\varrho^{m+1},
\end{aligned}}

worin $e_1, e_2, \ldots, e_s$ Konstanten sind, die nicht alle verschwinden,
und $\sigma_1, \sigma_2, \ldots, \sigma_s$ Funktionen, die in $\mathfrak{P}$ endlich sind;
$m$ ist die niedrigste Ordnungszahl in $\mathfrak{P}$, die unter den Funktionen $\eta_i$ vorkommt.

Der Inbegriff der Funktionen

\weq{3}{\eta = c_1\eta_1 + c_2\eta_2 + \cdots + c_s\eta_s}

mit Konstanten $c_1, c_2, \ldots, c_s$ heißt eine \gesperrt{Schar}. Wir bezeichnen sie durch

\weq{4}{(\eta_1, \eta_2, \ldots, \eta_s),}

und nennen $\eta_1, \eta_2, \ldots, \eta_s$ eine \gesperrt{Basis der Schar}. Die Ordnung
von $\eta$ in $\mathfrak{P}$ ist nicht kleiner als $m$ (sie ist größer, wenn
$e_1c_1 + e_2c_2 + \cdots + e_sc_s = 0$ ist). Es kann aber $\eta$ in keinem
Punkt $\mathfrak{P}$ unendlich werden, der nicht in $\mathfrak{A}$ enthalten ist, und
auch nicht in höhere Ordnung, als der Punkt $\mathfrak{P}$ in $\mathfrak{A}$ eingeht.

Wir können also setzen:

\weq{5}{\eta = \frac{\mathfrak{A}'}{\mathfrak{A}},}

worin $\mathfrak{A}'$ gleichfalls zur Klasse $A$ gehört.

Nimmt man statt $\mathfrak{A}$ ein anderes Polygon $\mathfrak{B}$ der Klasse $A$
und setzt

\weq{6}{\xi = \frac{\mathfrak{A}}{\mathfrak{B}},}

so folgt aus 1.:

\weq{7}{\begin{aligned}
\xi\eta_1 &= \eta_1' = \frac{\mathfrak{A}_1}{\mathfrak{B}},\\
\xi\eta_2 &= \eta_2' = \frac{\mathfrak{A}_2}{\mathfrak{B}},\\
&\text{\quad$\cdot\ \cdot\ \cdot\ \cdot\ \cdot\ \cdot\ \cdot\ \cdot$}\\
\xi\eta_s &= \eta_s' = \frac{\mathfrak{A}_s}{\mathfrak{B}}
\end{aligned}}

und

\weq{8}{\xi\eta = \eta' = c_1\eta_1' + c_2\eta_2' + \cdots + c_s\eta_s' = \frac{\mathfrak{A}'}{\mathfrak{B}}.}

Jedes durch den Nenner $\mathfrak{A}$ und ein Konstantensystem $c_1,
c_2, \ldots, c_s$ erzeugte Polygon $\mathfrak{A}'$ wird also auch durch jeden mit
%CONT>>
